% TOP_PROBLEM: {"id": 30006872, "title": "Koebe circle-domain uniformization conjecture", "category_id": 8, "category": {"id": 8, "name": "analysis", "display_name": "Analysis", "description": "Limits, continuity, calculus, and function theory.", "slug": "analysis", "order_index": 8, "created_at": "2026-07-31T15:26:25.671Z"}, "status": "open"}
% FIRST_POSTED: 2026-09-27
% !TeX program = lualatex
% ATTEMPT_STATUS: unresolved
\documentclass[11pt]{article}
\usepackage[margin=25mm]{geometry}
\usepackage{fontspec}
\setmainfont{Latin Modern Roman}
\usepackage{amsmath,amssymb,mathtools}
\usepackage{unicode-math}
\setmathfont{Latin Modern Math}
\usepackage{xurl,hyperref}
\hypersetup{colorlinks=true,urlcolor=blue}
\setcounter{secnumdepth}{0}
\setlength{\parindent}{0pt}
\setlength{\parskip}{5pt}
\setlength{\emergencystretch}{3em}
\begin{document}
\section*{\texorpdfstring{Koebe circle-domain uniformization conjecture}{Koebe circle-domain uniformization conjecture}}\label{30006872}
\subsection{Definitions and mathematical statement}
Write $\widehat{{\mathbb{C}}}$ for the Riemann sphere. A circle domain is a nonempty connected open subset $D\subseteq\widehat{{\mathbb{C}}}$ such that each connected component of $\widehat{{\mathbb{C}}}\setminus D$ is a closed round disk or a single point. The conjecture asserts that for every nonempty connected open $\Omega\subseteq\widehat{{\mathbb{C}}}$ there is a conformal bijection
\[
 f:\Omega\longrightarrow D
\]
onto a circle domain. A conformal bijection is a homeomorphism with a holomorphic inverse in local complex charts. The complement may have uncountably many components, including points; the sphere and the sphere minus a point are allowed. No uniqueness up to M\"obius transformations, continuous boundary extension, constructive procedure, or distortion bound is appended. A noninjective covering map is not the required uniformization.
\par\smallskip\textit{Formulation and concept references:} \href{https://www.math.stonybrook.edu/~bishop/classes/math536.S24/He_Schramm_1993.pdf}{[S1]}.

\subsection{Short English statement}
Is every domain on the Riemann sphere conformally equivalent to one whose complementary components are round disks or points, even with arbitrary connectivity?
\subsection{Research attempt: annular uniformization and the degeneration of circle complements}
\paragraph{1. Scope and source audit.}
The requested map is a conformal bijection onto a circle domain. A holomorphic covering map or a quasiconformal homeomorphism is not sufficient. The domain may have uncountably many complementary components, including point components. The 1993 He--Schramm source [S1] concerns the countably connected theorem. Ntalampekos's March 2026 survey [E1], Section 2, was consulted on 27 September 2026 for the formulation, partial results, and limitations of exhaustion approaches. It does not supply the unrestricted conjecture as a theorem. The discussion here proves explicit subcases and a precise failure of a naive limiting argument. Classical existence for the Dirichlet problem on a smooth bounded domain, the Riemann mapping theorem, and elementary holomorphic mapping theory are stated inputs.

\paragraph{2. Degenerate complements are part of the question.}
The whole sphere is itself a circle domain, with empty complement. The sphere minus one point is also a circle domain; stereographic projection identifies it with the complex plane, and a M\"obius map can move the missing point to any desired location. The punctured plane has two singleton complementary components on the sphere and is again already a circle domain. These cases must not be excluded by automatically normalizing every source domain into the unit disk.

If a simply connected domain is a proper subset of the plane and is not the whole plane, the Riemann mapping theorem gives a conformal bijection to the disk. The disk is a circle domain on the sphere because its complement is a closed round disk containing infinity. This proves a familiar subcase without imposing boundary regularity or asserting that the conformal map extends homeomorphically to the boundary.

\paragraph{3. A constructive doubly connected subcase.}
Let $\Omega$ be a bounded planar domain whose boundary consists of two disjoint smooth Jordan curves, an inner curve $\Gamma_0$ and an outer curve $\Gamma_1$. Solve the Dirichlet problem for a harmonic function $u$ with $u=0$ on $\Gamma_0$ and $u=1$ on $\Gamma_1$. The maximum principle gives $0<u<1$ in the domain. Define its capacity
\[
 C=\int_\Omega |\nabla u|^2\,dA>0.
\]
Green's identity shows that the outward flux through the outer curve is $C$; the flux through the inner curve is $-C$. The closed one-form $*du=-u_y\,dx+u_x\,dy$ has period $C$ around a counterclockwise loop surrounding the hole. A local primitive $v$ therefore changes by an integer multiple of $C$ under continuation around any closed loop. This last statement uses that the fundamental group of a Jordan annulus is infinite cyclic.

The function
\[
 F(z)=\exp\left(\frac{2\pi}{C}(u(z)+iv(z))\right)
\]
is consequently single-valued and holomorphic. It maps $\Omega$ into the round annulus $1<|w|<\exp(2\pi/C)$ and extends continuously to the two boundary curves with their respective boundary radii. It is proper: the inverse image of a compact set staying away from the two target circles stays away from both source boundaries. Its winding number around the origin on a positively oriented generator is one, because $v$ increases by $C$.

A nonconstant proper holomorphic map between these annuli has a well-defined positive degree. That degree equals the induced winding number on a generator, hence equals one. Equivalently, apply the argument principle to $F-w$ using the two boundary curves: the outer image winds once around any point in the target annulus and the inner image winds zero times, so each target point has exactly one preimage counted with multiplicity. It follows that $F$ is bijective and has no critical point. The inverse is holomorphic by the inverse function theorem. Thus the desired conformal bijection is obtained explicitly from the harmonic potential.

There is no hidden assumption that $\nabla u$ is nonzero everywhere. Degree one proves that afterward, through $F'=(2\pi/C)F(u_x-iu_y)$. The existence and boundary regularity of $u$ are the standard analytic inputs appropriate to the smooth Jordan boundaries; arbitrary infinitely connected boundaries are not covered by this construction.

\paragraph{4. Capacity normalization checked directly on round annuli.}
For $A(r,R)=\{r<|z|<R\}$ with $0<r<R<\infty$, the normalized potential is
\[
 u(z)=\frac{\log(|z|/r)}{\log(R/r)},\qquad
 C=\frac{2\pi}{\log(R/r)}.
\]
Indeed $|\nabla u|=1/(|z|\log(R/r))$, and polar integration gives the displayed energy. This is the minimum energy among functions taking the prescribed boundary values. For each angle $\theta$, Cauchy--Schwarz yields
\[
 1=\left(\int_r^R\partial_\rho f(\rho,\theta)\,d\rho\right)^2
 \leq\left(\int_r^R|\partial_\rho f|^2\rho\,d\rho\right)
       \left(\int_r^R\frac{d\rho}{\rho}\right).
\]
Integrating in $\theta$ and dropping the nonnegative angular part of the Dirichlet energy proves the lower bound $2\pi/\log(R/r)$. Equality holds for $u$. In particular the construction in Section 3 gives outer-to-inner radius ratio $\exp(2\pi/C)$, not its reciprocal.

In two dimensions the Dirichlet integral is invariant under conformal changes of coordinates: the squared derivative scaling cancels the Jacobian. Thus $\log(R/r)/(2\pi)$ is a conformal invariant of an annulus. Letting $r$ tend to zero gives a different limiting type, a punctured disk, with zero capacity for the puncture. Setting $C=0$ directly in the exponential formula is invalid; the degenerate case requires its own treatment.

\paragraph{5. Uncountability by itself is not an obstruction.}
Let $E$ be the usual ternary Cantor set placed in a bounded interval of the real axis, and put $\Omega=\widehat{\mathbb C}\setminus E$. The set $\Omega$ is open and connected. To check connectedness concretely, points off the real axis can be connected in an upper or lower half-plane; the two half-planes can be joined around an endpoint of the containing interval. A real point outside $E$ can first be moved a short vertical distance through a small disk disjoint from $E$. These paths connect all points in the complement.

Every connected component of $E$ is a singleton. Consequently $\Omega$ is already a circle domain with uncountably many point complementary components. Thus the conjecture cannot be reduced to countably connected domains by asserting that an arbitrary circle domain has only countably many components. What is countable is the family of its nondegenerate complementary disks: each has positive spherical area, and pairwise disjoint positive-area subsets of a finite measure space form a countable family. Singleton components are not counted by that measure argument.

\paragraph{6. Round complementary components need not survive a kernel limit.}
Here is a finite, explicit degeneration. For each integer $j\geq1$, let the centers be $c_{j,k}=-1+2k/j$, $0\leq k\leq j$, and let the common radius be $r_j=1/(10j^2)$. The disks $\overline B(c_{j,k},r_j)$ are pairwise disjoint, because their center spacing is $2/j>2r_j$. Define
\[
 D_j=\widehat{\mathbb C}\setminus
 \bigcup_{k=0}^{j}\overline B(c_{j,k},r_j).
\]
Each $D_j$ is a circle domain containing infinity. Their kernel with base point infinity is
\[
 D=\widehat{\mathbb C}\setminus[-1,1].
\]
To prove this, every compact subset of $D$ has positive distance from the segment and is eventually disjoint from all deleted disks. On the other hand, every neighborhood of every point of $[-1,1]$ contains one of the centers for all sufficiently large $j$, since the mesh tends to zero. No such neighborhood is eventually contained in $D_j$. These two statements characterize the kernel, whose indicated complement is connected.

The kernel $D$ is not a circle domain: its one complementary component is a nondegenerate segment, which is neither a closed round disk nor a point. Of course $D$ itself can be conformally uniformized; this example is not a counterexample to Koebe's conjecture. It proves exactly that the geometric property ``all complementary components are round disks or points'' is not closed under this natural limit of domains. A limiting proof needs additional control preventing components from coalescing into a nonround continuum.

\paragraph{7. Why normal families do not complete the argument.}
If univalent holomorphic maps on an increasing exhaustion have a locally uniform nonconstant limit, the limit is univalent. One proof applies Hurwitz's theorem to the differences $f_j(z)-f_j(w)$ for fixed $w$, after removing the forced zero at $z=w$. A normalization ensuring a nonzero limiting derivative prevents a constant limit when the family is locally bounded. These are valuable analytic steps, but neither one says that every complementary component of the image remains round. Section 6 isolates the missing geometric compactness statement.

There is also a separate normalization issue: mapping all exhaustion domains into one fixed bounded target with a uniform lower derivative bound is a condition to prove, not a free consequence of choosing three labels. And the universal cover of an annulus is simply connected, but the covering projection is not injective. Uniformizing that cover solves a different mapping problem.

\paragraph{8. Diagnostics and conclusion.}
The saved exact checks verify disjointness of the rational disks in Section 6 and the mesh/radius estimates for a finite range. The kernel statement for all $j$ follows from the displayed inequalities and the compactness argument. The capacity identity is analytic and is not certified by sampling a few annuli.

\textbf{UNRESOLVED.} The attempt constructs the required map for smooth doubly connected domains and checks basic degenerate and Cantor-complement examples. The general gap is to control nonround limiting complementary continua while obtaining a nondegenerate conformal limit for an arbitrary domain. Neither countability of positive-area disks nor normal-family convergence supplies that control.

\subsection{Sources}
\begin{itemize}
\item[S1]Zheng-Xu He and Oded Schramm, Fixed Points, Koebe Uniformization and Circle Packings, Annals of Mathematics137(1993),369-406; with Ntalampekos2026 survey.
\url{https://www.math.stonybrook.edu/~bishop/classes/math536.S24/He_Schramm_1993.pdf}.
\textit{Formulation source.}
\item[E1]D. Ntalampekos, Uniformization problems in the plane: A survey, arXiv:2603.15098v1, 16 March 2026, Section 2.
\url{https://arxiv.org/html/2603.15098v1}.
\textit{Primary author survey consulted 27 September 2026 for the unrestricted conjecture and the distinction between existence, rigidity, and partial countability results.}
\end{itemize}
\end{document}
